
\scanpage{59}
\mspage{15}
pour \deleted{traduire} traduire le groupoïde de Teichmüller topologique, on ne trouvera jamais que des structures où les cercles de découpage sont effectivement marqués.
\deleted{Si on désigne désormais par $S\mathcal TD_{\mathrm{top}}^\pm$ le groupoïde}
\inserted{Néanmoins on} veut avoir le droit de les ``gommer'', de les oublier \inserted{\unsure{les|ces} marquages} --- même s'ils y sont bel et bien!

D) \uline{Opération de recollement}. C'est la plus cruciale, la moins tautologique peut-être de toutes les opérations que nous passons en revue.
Pour la pratiquer, il faut comme donnée préliminaire

a) Une partie \deleted{\unsure{les}} $A'$ de $A_l$ (arêtes libres), et une involution ss pts fixes (i.e. une relation d'équivalence à classes de cardinal 2) $\sigma'$ sur $A'$.
\deleted{\unsure{et les}\ $\unsure{\sigma}$}

b) Pour \deleted{$\{a,a'\}\in\unsure{A'_l}$} \unsure{toute} orbite \deleted{de} $\{a,a'\}$ de $A'$, un isom. renversant l'orientation de $C_a$ sur $C_{a'}$ et de $C_{a'}$ sur $C_a$, inverses l'un de l'autre.
Pour que l'opération soit bien déterminée à isom unique
\mspage{58}
